% GENERATED FILE. Edit canonical lessons, then run npm run generate:books.
% canonical-source-hash: fnv1a64:b953596bbb7982bf
% canonical-lessons: JA-C421-negiru, JA-C421-teika, JA-C421-neage, JA-C421-nesage, JA-C421-shouhin

\chapter{Haggling at the Market}
\label{ch:ja-haggling-at-the-market}
\hlchaptermodality{421}

\begin{chapteropening}
\textbf{By the end of this chapter:} \emph{I can talk about prices going up and down.}

Complete the last lesson of chapter 421: I can talk about prices going up and down.
\end{chapteropening}

\section[\texorpdfstring{\ja{ねぎる} (negiru)}{ねぎる (negiru)}]{\ja{ねぎる} (negiru) — to haggle}

\label{lesson:JA-C421-negiru}



\begin{quote}
Before the new one: say the Japanese for to pay a bill, then the Japanese for to refund.
\end{quote}

\subsection*{You'll want to know: \ja{ねぎる}}
\textbf{\ja{ねぎる}} — \emph{negiru} — \textquotedblleft{}to haggle\textquotedblright{}.

\begin{quote}
\textbf{\ja{いちばで} \ja{ねぎって、やすく} \ja{かいました。}} — \emph{ichiba de negitte, yasuku kaimashita} — \textquotedblleft{}I haggled at the market and bought it cheaply.\textquotedblright{}
\end{quote}

Polite \textbf{\ja{ねぎります}}, \emph{negirimasu}; a godan verb. Past (polite): \textbf{\ja{ははが} \ja{いちばで} \ja{ねぎりました。}}, \emph{haha ga ichiba de negirimashita}.

\begin{cousinweb}
\textbf{\ja{ねぎる}} (\emph{negiru}) is a native Japanese word (wago), inherited rather than borrowed from Chinese, built from \emph{ne}, \textquotedblleft{}price\textquotedblright{} and \emph{giru}, \textquotedblleft{}cutting\textquotedblright{} (from \emph{kiru}); its kanji were attached later as kun'yomi (JMdict).
\end{cousinweb}

\begin{grammarlens}[title={What you've built}]
One more word for getting around.
\end{grammarlens}

\subsection*{Your turn}
\begin{itemize}
\raggedright
  \item \emph{Say it:} \emph{negiru}
  \item \emph{Say it:} \emph{negiru}, once more
  \item \emph{Say it:} \emph{ichiba de negitte, yasuku kaimashita} — I haggled at the market and bought it cheaply.
  \item \emph{Recall:} say the Japanese for to pay a bill, then the Japanese for to refund, then say \emph{negiru} again
\end{itemize}

\begin{tcolorbox}[breakable,colback=teal!4,colframe=teal!35!black,title={Before you move on}]
What does \ja{ねぎる} mean? (\textquotedblleft{}To haggle\textquotedblright{}.) Say it once more.
\end{tcolorbox}
\section[\texorpdfstring{\ja{ていか} (teika)}{ていか (teika)}]{\ja{ていか} (teika) — the list price}

\label{lesson:JA-C421-teika}



\begin{quote}
Before the new one: say the Japanese for to refund, then the Japanese for to haggle.
\end{quote}

\subsection*{You'll want to know: \ja{ていか}}
\textbf{\ja{ていか}} — \emph{teika} — \textquotedblleft{}the list price\textquotedblright{}.

\begin{quote}
\textbf{\ja{わりびきで} \ja{ていかより} \ja{やすく} \ja{かいました。}} — \emph{waribiki de teika yori yasuku kaimashita} — \textquotedblleft{}With the discount I bought it below the list price.\textquotedblright{}
\end{quote}

\begin{cousinweb}
\textbf{\ja{ていか}} (\emph{teika}) is Sino-Japanese (kango): two Chinese characters read with their on'yomi, readings borrowed from Chinese. They are \emph{tei}, \textquotedblleft{}fix\textquotedblright{} (Middle Chinese \emph{dengH}) and \emph{ka}, \textquotedblleft{}value\textquotedblright{} (\emph{kaeH}) (KANJIDIC2; Middle Chinese in Baxter's transcription).
\end{cousinweb}

\begin{grammarlens}[title={What you've built}]
One more word for getting around.
\end{grammarlens}

\subsection*{Your turn}
\begin{itemize}
\raggedright
  \item \emph{Say it:} \emph{teika}
  \item \emph{Say it:} \emph{teika}, once more
  \item \emph{Say it:} \emph{waribiki de teika yori yasuku kaimashita} — With the discount I bought it below the list price.
  \item \emph{Recall:} say the Japanese for to refund, then the Japanese for to haggle, then say \emph{teika} again
\end{itemize}

\begin{tcolorbox}[breakable,colback=teal!4,colframe=teal!35!black,title={Before you move on}]
What does \ja{ていか} mean? (\textquotedblleft{}The list price\textquotedblright{}.) Say it once more.
\end{tcolorbox}
\section[\texorpdfstring{\ja{ねあげ} (neage)}{ねあげ (neage)}]{\ja{ねあげ} (neage) — a price rise}

\label{lesson:JA-C421-neage}



\begin{quote}
Before the new one: say the Japanese for to haggle, then the Japanese for the list price.
\end{quote}

\subsection*{You'll want to know: \ja{ねあげ}}
\textbf{\ja{ねあげ}} — \emph{neage} — \textquotedblleft{}a price rise\textquotedblright{}.

\begin{quote}
\textbf{\ja{らいげつ} \ja{うんちんの} \ja{ねあげが} \ja{あります。}} — \emph{raigetsu unchin no neage ga arimasu} — \textquotedblleft{}Fares are going up next month.\textquotedblright{}
\end{quote}

\begin{cousinweb}
\textbf{\ja{ねあげ}} (\emph{neage}) is a native Japanese word (wago), inherited rather than borrowed from Chinese, built from \emph{ne}, \textquotedblleft{}price\textquotedblright{} and \emph{age}, \textquotedblleft{}raising\textquotedblright{} (from \emph{ageru}); its kanji were attached later as kun'yomi (JMdict).
\end{cousinweb}

\begin{grammarlens}[title={What you've built}]
One more word for getting around.
\end{grammarlens}

\subsection*{Your turn}
\begin{itemize}
\raggedright
  \item \emph{Say it:} \emph{neage}
  \item \emph{Say it:} \emph{neage}, once more
  \item \emph{Say it:} \emph{raigetsu unchin no neage ga arimasu} — Fares are going up next month.
  \item \emph{Recall:} say the Japanese for to haggle, then the Japanese for the list price, then say \emph{neage} again
\end{itemize}

\begin{tcolorbox}[breakable,colback=teal!4,colframe=teal!35!black,title={Before you move on}]
What does \ja{ねあげ} mean? (\textquotedblleft{}A price rise\textquotedblright{}.) Say it once more.
\end{tcolorbox}
\section[\texorpdfstring{\ja{ねさげ} (nesage)}{ねさげ (nesage)}]{\ja{ねさげ} (nesage) — a price cut}

\label{lesson:JA-C421-nesage}



\begin{quote}
Before the new one: say the Japanese for the list price, then the Japanese for a price rise.
\end{quote}

\subsection*{You'll want to know: \ja{ねさげ}}
\textbf{\ja{ねさげ}} — \emph{nesage} — \textquotedblleft{}a price cut\textquotedblright{}.

\begin{quote}
\textbf{\ja{この} \ja{みせは} \ja{ゆうがた} \ja{ねさげを} \ja{します。}} — \emph{kono mise wa yūgata nesage o shimasu} — \textquotedblleft{}This shop cuts its prices in the evening.\textquotedblright{}
\end{quote}

\begin{cousinweb}
\textbf{\ja{ねさげ}} (\emph{nesage}) is a native Japanese word (wago), inherited rather than borrowed from Chinese, built from \emph{ne}, \textquotedblleft{}price\textquotedblright{} and \emph{sage}, \textquotedblleft{}lowering\textquotedblright{} (from \emph{sageru}); its kanji were attached later as kun'yomi (JMdict).
\end{cousinweb}

\begin{grammarlens}[title={What you've built}]
One more word for getting around.
\end{grammarlens}

\subsection*{Your turn}
\begin{itemize}
\raggedright
  \item \emph{Say it:} \emph{nesage}
  \item \emph{Say it:} \emph{nesage}, once more
  \item \emph{Say it:} \emph{kono mise wa yūgata nesage o shimasu} — This shop cuts its prices in the evening.
  \item \emph{Recall:} say the Japanese for the list price, then the Japanese for a price rise, then say \emph{nesage} again
\end{itemize}

\begin{tcolorbox}[breakable,colback=teal!4,colframe=teal!35!black,title={Before you move on}]
What does \ja{ねさげ} mean? (\textquotedblleft{}A price cut\textquotedblright{}.) Say it once more.
\end{tcolorbox}
\section[\texorpdfstring{\ja{しょうひん} (shōhin)}{しょうひん (shōhin)}]{\ja{しょうひん} (shōhin) — merchandise; a product}

\label{lesson:JA-C421-shouhin}



\begin{quote}
Before the new one: say the Japanese for a price rise, then the Japanese for a price cut.
\end{quote}

\subsection*{You'll want to know: \ja{しょうひん}}
\textbf{\ja{しょうひん}} — \emph{shōhin} — \textquotedblleft{}merchandise; a product\textquotedblright{}.

\begin{quote}
\textbf{\ja{この} \ja{みせの} \ja{しょうひんは} \ja{たかいです。}} — \emph{kono mise no shōhin wa takai desu} — \textquotedblleft{}The goods in this shop are expensive.\textquotedblright{}
\end{quote}

\begin{cousinweb}
\textbf{\ja{しょうひん}} (\emph{shōhin}) is Sino-Japanese (kango): two Chinese characters read with their on'yomi, readings borrowed from Chinese. They are \emph{shō}, \textquotedblleft{}selling\textquotedblright{} (Middle Chinese \emph{syang}) and \emph{hin}, \textquotedblleft{}goods\textquotedblright{} (\emph{phimX}) (KANJIDIC2; Middle Chinese in Baxter's transcription).
\end{cousinweb}

\begin{grammarlens}[title={What you've built}]
One more word for getting around.
\end{grammarlens}

\subsection*{Your turn}
\begin{itemize}
\raggedright
  \item \emph{Say it:} \emph{shōhin}
  \item \emph{Say it:} \emph{shōhin}, once more
  \item \emph{Say it:} \emph{kono mise no shōhin wa takai desu} — The goods in this shop are expensive.
  \item \emph{Recall:} say the Japanese for a price rise, then the Japanese for a price cut, then say \emph{shōhin} again
\end{itemize}

\begin{tcolorbox}[breakable,colback=teal!4,colframe=teal!35!black,title={Before you move on}]
What does \ja{しょうひん} mean? (\textquotedblleft{}Merchandise; a product\textquotedblright{}.) Say it once more.
\end{tcolorbox}
